\begin{abstract}
% \todo{"Favourably" might not be well evidenced.}
Large language models exhibit recurring behavioural patterns, often described as personas. These personas have been found to be useful for controlling generalisation and improving safety, but we lack reliable tools for decomposing, measuring, and controlling these patterns. Our central insight is to treat personas as positions in a space of established behavioural traits, using the OCEAN framework to describe model personas in terms of \textbf{Openness}, \textbf{Conscientiousness}, \textbf{Extroversion}, \textbf{Agreeableness}, and \textbf{Neuroticism}. We study the variation of personas as movement along axes in weight space by training low-rank adapters to amplify or suppress individual traits, and evaluating their effects using calibrated LLM judges, trait-specific multiple-choice benchmarks, and standard capability evaluations. We find that these adapters induce predictable trait shifts: scaling an adapter produces mostly monotonic behavioural change over useful ranges, and linear combinations of adapters approximately recover the effects of their components. These interventions preserve performance on standard capability benchmarks at moderate scales. We further show that the induced trait axes affect safety-relevant behaviour in downstream evaluations: for example, moving along neuroticism and agreeableness axes affects frustration and sycophancy respectively. We also introduce an unsupervised psychometric pipeline that discovers interpretable latent behavioural factors from model rollouts. Our results show that LLM personas can be analyzed through a structured ‘trait space’ and controlled using a partially-composable weight space, providing a bridge between personality measurement, model editing, and safety-relevant behavioural control.
\end{abstract}
